\section{Existence of Pfaffian point processes} \label{sect:existence_PfPP}
\subsection{Fock representations}
Here, we see Fock representations of $\mcal{C}(\mcal{K},\Gamma)$, which play prominent roles in the theory of a CAR algebra.
\subsubsection{General construction}
Let $\mcal{H}$ be a Hilbert space.
For each $n\in\mbb{N}$, we write $\bigwedge^{n}\mcal{H}$ for the $n$-th wedge product of $\mcal{H}$, which is generated by the vectors $u_{1}\wedge\cdots \wedge u_{n}$, $u_{1},\dots, u_{n}\in\mcal{H}$ subject to the anti-symmetry:
\begin{equation*}
	u_{1}\wedge\cdots \wedge u_{n}=\operatorname{sgn} (\sigma) u_{\sigma(1)}\wedge\cdots \wedge u_{\sigma (n)},\qquad \sigma\in \mfrak{S}_{n}.
\end{equation*}
When we take $\{e_{i}\}_{i=1,2,\dots}$ for a complete orthonormal system of $\mcal{H}$, then the vectors $e_{i_{1}}\wedge\cdots \wedge e_{i_{n}}$, $i_{1}>\cdots >i_{n}$ form a complete orthonormal system of $\bigwedge^{n}\mcal{H}$.
The Fermi Fock space over $\mcal{H}$ is defined by
\begin{equation*}
	\mcal{F}(\mcal{H})=\bigoplus_{n=0}^{\infty}\bigwedge^{n}\mcal{H},
\end{equation*}
where we set $\bigwedge^{0}\mcal{H}=\mbb{C}\bm{1}$.
The Fermi Fock space admits the natural inner product induced from each component of the direct sum and becomes a Hilbert space, i.e., the direct sum is understood in the topological sense.

For each $u\in \mcal{H}$, the creation operator $a_{\mcal{H}}^{\ast}(u)$ is an operator on $\mcal{F}(\mcal{H})$ defined by
\begin{equation*}
	a_{\mcal{H}}^{\ast}(u)\colon \ \eta\mapsto u\wedge \eta,\qquad \eta\in \mcal{F}(\mcal{H}).
\end{equation*}
The annihilation operator $a_{\mcal{H}}(u)$ is defined as the adjoint operator of $a_{\mcal{H}}^{\ast}(u)$.
By definition, it is obvious that creation and annihilation operators act as
\begin{equation*}
	a_{\mcal{H}}^{\ast}(u)\colon \ \bigwedge^{n}\mcal{H}\to \bigwedge^{n+1}\mcal{H},\qquad a_{\mcal{H}}(u)\colon \ \bigwedge^{n}\mcal{H}\to \bigwedge^{n-1}\mcal{H}.
\end{equation*}
We can also see that the assignment $u\mapsto a_{\mcal{H}}^{\ast}(u)$ is linear and $u\mapsto a_{\mcal{H}}(u)$ is anti-linear.

\subsubsection{Fock representation and a quasi-free state}
Let us take a projection operator $P\in\mrm{Gr}(\mcal{K},\Gamma)$, associated to which we can construct a representation of $\mcal{C}(\mcal{K},\Gamma)$ on a Fock space $\mcal{F}(P\mcal{K})$.
To describe the action, notice that, from the property $1-P=\overline{P}$, we can see that the projection to the complementary subspace of $P\mcal{K}$ is $\overline{P}$.
Let us set
\begin{equation*}
	\pi_{P}\colon \ \mcal{C}(\mcal{K},\Gamma)\to \mfrak{B}(\mcal{F}(P\mcal{K}));\qquad B(f)\mapsto a_{P\mcal{K}}^{\ast}(Pf)+a_{P\mcal{K}}(P\Gamma f),\qquad f\in\mcal{K}.
\end{equation*}
Then, it is known that $(\pi_{P},\mcal{F}(P\mcal{K}))$ is a faithful and irreducible representation of $\mcal{C}(\mcal{K},\Gamma)$.
When we set
\begin{equation*}
	\varphi_{P}(A)=(\bm{1}, \pi_{P}(A)\bm{1})_{\mcal{F}(P\mcal{K})},\qquad A\in \mcal{C}(\mcal{K},\Gamma),
\end{equation*}
we can verify that $\varphi_{P}$ is exactly the quasi-free state corresponding to $P$ in the sense that it possesses the property required in Lemma \ref{lem:quasi-free_covariance}

In particular, when we take $P_{0}\in\mrm{Gr}(\mcal{K},\Gamma)$, we have $P_{0}\mcal{K}\simeq \ell^{2}(\mfrak{X})$, and the map $\pi_{P_{0}}$ is described as
\begin{equation*}
	\pi_{P_{0}}(a^{\ast}(u))=a_{P_{0}\mcal{K}}^{\ast}(u),\qquad \pi_{P_{0}}(a(u))=a_{P_{0}\mcal{K}}(u),\qquad u\in P_{0}\mcal{K}\simeq \ell^{2}(\mfrak{X}).
\end{equation*}

Let us describe a standard complete orthonormal system of the Fock space $\mcal{F}\big(\ell^{2}(\mfrak{X})\big)$. Since~$\mfrak{X}$ is countable, it can be equipped with a linear order~$\le$. For each $\omega=\{x_{1}>\cdots>x_{n}\}\in\Omega^{\circ}$, we set
\begin{equation*}
	e_{\omega}:=e_{x_{1}}\wedge\cdots \wedge e_{x_{n}}.
\end{equation*}
Then the collection $\{e_{\omega}\}_{\omega\in\Omega^{\circ}}$ forms a complete orthonormal system.

\begin{proof}[Proof of Lemma \ref{lem:quasi-free_covariance}]
Let $\varphi$ be a quasi-free state over $\mcal{C}(\mcal{K},\Gamma)$.
Then the assignment
\begin{equation*}
	Q_{\varphi}\colon \ \mcal{K}\times \mcal{K}\to\mbb{C};\qquad (f,g)\mapsto \varphi (B(f)^{\ast}B(g))
\end{equation*}
defines a quadratic form on $\mcal{K}$.
It follows from the relation in the CAR algebra that
\begin{equation*}
	B(f)^{\ast}B(f)\le B(f)^{\ast}B(f)+B(f)B(f)^{\ast}=\|f\|^{2},\qquad f\in\mcal{K}.
\end{equation*}
Since the norm of a state is unity, we have $Q_{\varphi}(f,f)\le \|B(f)^{\ast}B(f)\|\le \|f\|^{2}$ for any $f\in\mcal{K}$, which implies that $Q_{\varphi}$ is a bounded quadratic form. Therefore, owing to the correspondence between bounded operators and bounded quadratic forms (see, e.g., \cite[Chapter~1]{Koshmanenko1999}), there exists a bounded operator $S\in \mfrak{B}(\mcal{K})$ such that $Q_{\varphi}(f,g)=(f,Sg)_{\mcal{K}}$, $f,g \in \mcal{K}$. It also follows from the positivity of the state that the quadratic form $Q_{\varphi}$ is positive, and therefore, $S=S^{\ast}\ge 0$. Again, from the relation in the CAR algebra, we have
\begin{equation*}
	(f,Sg)_{\mcal{K}}=(f,g)_{\mcal{K}}-(\Gamma g, S\Gamma f)_{\mcal{K}}=(f,g)_{\mcal{K}}-\big(f,\overline{S} g\big)_{\mcal{K}},\qquad f,g\in\mcal{K},
\end{equation*}
which implies that $\Gamma S \Gamma=1-S$. Since $\Gamma S \Gamma\ge 0$, we can see that $S\le 1$.

Conversely, given an operator $S\in\mcal{Q}(\mcal{K},\Gamma)$, we define the following operator
\begin{equation*}
	P_{S}=\left(
	\begin{matrix}
	S & S^{1/2}(1-S)^{1/2} \\
	S^{1/2}(1-S)^{1/2} & 1-S
	\end{matrix}
	\right)
\end{equation*}
acting on $\what{\mcal{K}}=\mcal{K}\oplus\mcal{K}$. Notice that the assumption $0\le S=S^{\ast}\le 1$ ensures that the square roots~$S^{1/2}$ and $(1-S)^{1/2}$ make sense. Let us equip this Hilbert space with the anti-unitary involution
\begin{equation*}
	\what{\Gamma}=\left(
	\begin{matrix}
	\Gamma & 0 \\
	0 & -\Gamma
	\end{matrix}
	\right).
\end{equation*}
Then, it can be checked that $P_{S}\in \mrm{Gr}\big(\what{\mcal{K}},\what{\Gamma}\big)$, which implies that the functional $\varphi_{P_{S}}$ defined by
\begin{equation*}
	\varphi_{P_{S}}(A)= (\bm{1},\pi_{P_{S}}(A)\bm{1} )_{\mcal{F}(P_{S}\what{\mcal{K}})},\qquad A\in \mcal{C}\big(\what{\mcal{K}},\what{\Gamma}\big)
\end{equation*}
is a quasi-free state over $\mcal{C}\big(\what{\mcal{K}},\what{\Gamma}\big)$. Now, since $\what{\Gamma}$ acts diagonally along the direct sum decomposition, $\mcal{C}(\mcal{K},\Gamma)$ is regarded as a subalgebra of $\mcal{C}\big(\what{\mcal{K}},\what{\Gamma}\big)$ and the restriction $\varphi_{S}=\varphi_{P_{S}}|_{\mcal{C}(\mcal{K},\Gamma)}$ is the quasi-free state corresponding to the given~$S$.
\end{proof}

\subsection{Proof of Proposition \ref{prop:existence_ppp}}
Now, we are at the position of proving Proposition~\ref{prop:existence_ppp}.
Our strategy is to check the criteria by Lenard \cite{Lenard1975a, Lenard1975b}.

\subsubsection{Lenard's criteria}
Suppose that a system of functions $\big\{\rho_{n}\colon\mfrak{X}^{n}\to\mbb{R}\big\}_{n=1}^{\infty}$ is given.
A question is if there exists a~probability measure~$M$ on $(\Omega,\Sigma)$ such that its $n$-point correlation function $\rho^{M}_{n}$ coincides with the given function $\rho_{n}$ for every $n\in\mbb{N}$.
Lenard \cite{Lenard1975a, Lenard1975b} clarified the necessary and sufficient conditions for this to happen.

\begin{Theorem}[\cite{Lenard1975a, Lenard1975b}]
A system of functions $\big\{\rho_{n}\colon \mfrak{X}^{n}\to\mbb{R}\big\}_{n=1}^{\infty}$ is a one of correlation functions for a probability measure on $(\Omega,\Sigma)$ if it possesses the following properties:
\begin{enumerate}\itemsep=0pt
\item[$1.$] {\it Symmetry:} each function $\rho_{n}$ is a symmetric function, i.e.,
	\begin{equation*}
		\rho_{n}(x_{\sigma (1)},\dots, x_{\sigma (n)})=\rho_{n}(x_{1},\dots, x_{n})
	\end{equation*}
	for arbitrary $\sigma\in\mfrak{S}_{n}$ and $x_{1},\dots, x_{n}\in\mfrak{X}$.
\item[$2.$] {\it Positivity:} for any system of functions $\bm{\Phi}=\big\{\Phi_{n}\colon \mfrak{X}^{n}\to\mbb{R}\big\}_{n=0}^{N}$ $(\Phi_{0}$ is understood as a~constant$)$ with finite support such that
	\begin{equation}\label{eq:positivity_test_function}
		\Phi_{0}+\sum_{n=1}^{N}\sum_{\substack{i_{1},\dots,i_{n}\in I \\ \mrm{distinct}}}\Phi_{n}(x_{i_{1}},\dots, x_{i_{n}})\ge 0
	\end{equation}
	holds for any $\omega=\{x_{i}\}_{i\in I}\in\Omega$, we have
	\begin{equation*}
		\Phi_{0}+\sum_{n=1}^{N}\sum_{x_{1},\dots, x_{n}\in\mfrak{X}}\Phi_{n}(x_{1},\dots, x_{n})\rho_{n}(x_{1},\dots, x_{n})\ge 0.
	\end{equation*}
\end{enumerate}
\end{Theorem}

Therefore, Proposition \ref{prop:existence_ppp} reduces to the following assertion.
\begin{Lemma}
\label{lem:existence_ppp}
Let $S\in\mcal{Q}(\mcal{K},\Gamma)$. Then the system of functions
\begin{equation*}
	\rho_{n}(x_{1},\dots, x_{n})=\pf \left[\mbb{K}_{S}(x_{i},x_{j})\right]_{1\le i,j\le n}, \qquad x_{1},\dots, x_{n}\in\mfrak{X}, \qquad n\in\mbb{N}
\end{equation*}
satisfies the symmetry and positivity conditions.
\end{Lemma}

\subsubsection{Proof of Lemma \ref{lem:existence_ppp}}
As we have seen, each correlation function is realized as
\begin{equation*}
	\rho_{n}(x_{1},\dots, x_{n})=\varphi_{S}\big(a_{x_{1}}^{\ast}\cdots a_{x_{n}}^{\ast}a_{x_{n}}\cdots a_{x_{1}}\big).
\end{equation*}
Then, the symmetry condition is obviously satisfied since
\begin{equation*}
	a_{x_{1}}^{\ast}\cdots a_{x_{n}}^{\ast}a_{x_{n}}\cdots a_{x_{1}}
	=a_{x_{\sigma (1)}}^{\ast}\cdots a_{x_{\sigma (n)}}^{\ast}a_{x_{\sigma (n)}}\cdots a_{x_{\sigma (1)}}
\end{equation*}
holds for all $\sigma\in\mfrak{S}_{n}$.

To show the positivity condition, let $\bm{\Phi}=\big\{\Phi_{n}\colon\mfrak{X}^{n}\to\mbb{R}\big\}_{n=0}^{N}$ be a system of functions satisfying~(\ref{eq:positivity_test_function}).
We can see that
\begin{equation*}
	\Phi_{0}+\sum_{n=1}^{N}\sum_{x_{1},\dots, x_{n}\in\mfrak{X}}\Phi_{n}(x_{1},\dots, x_{n})\rho_{n}(x_{1},\dots, x_{n})
	=\varphi_{S} (A_{\bm{\Phi}} ),
\end{equation*}
where we set
\begin{equation*}
	A_{\bm{\Phi}}=\Phi_{0}+\sum_{n=1}^{N}\sum_{x_{1},\dots,x_{n}\in\mfrak{X}}\Phi_{n}(x_{1},\dots,x_{n})a_{x_{1}}^{\ast}\cdots a_{x_{n}}^{\ast}a_{x_{n}}\cdots a_{x_{1}}.
\end{equation*}
Therefore, owing to the positivity of a state over a C$^{\ast}$-algebra, it suffices to show that $A_{\bm{\Phi}}\in\mcal{C}(\mcal{K},\Gamma)$ is a positive element, which can be checked in any faithful representation. In fact, when we take a faithful representation $(\pi,\mcal{H})$, we may regard $\mcal{C}(\mcal{K},\Gamma)$ as a C$^{\ast}$-subalgebra of $\mfrak{B}(\mcal{H})$ via the embedding~$\pi$. Since the spectrum of an element in a C$^{\ast}$-subalgebra coincides with that in a~whole algebra (see, e.g., \cite[Proposition~2.2.7]{BratteliRobinson1987}), the relevant element $A_{\bm{\Phi}}\in\mcal{C}(\mcal{K},\Gamma)$ is positive if and only if $\pi \left(A_{\bm{\Phi}}\right)$ is a positive operator on $\mcal{H}$.

We can, in particular, take a Fock representation $\big(\pi_{P_{0}},\mcal{F}\big(\ell^{2}(\mfrak{X})\big)\big)$.
\begin{Lemma}
In the Fock space $\mcal{F}\big(\ell^{2}(\mfrak{X})\big)$, the complete orthonormal system $\{e_{\omega}\}_{\omega\in\Omega^{\circ}}$ diagonalizes $\pi_{P_{0}} (A_{\bm{\Phi}} )$ so that
\begin{equation*}
	\pi_{P_{0}}\left(A_{\bm{\Phi}}\right) e_{\omega}=\left(\Phi_{0}+\sum_{k=1}^{N}\sum_{\substack{i_{1},\dots,i_{k}=1\\ \mrm{distinct}}}^{n}\Phi_{k}(x_{i_{1}},\dots, x_{i_{k}})\right) e_{\omega},\qquad \omega=\{x_{1}>\cdots >x_{n}\},
\end{equation*}
\end{Lemma}
\begin{proof}
This follows from a direct computation.
Let us notice that $a_{x_{1}}^{\ast}\cdots a_{x_{n}}^{\ast}a_{x_{n}}\cdots a_{x_{1}}=0$ if any two points from $x_{1},\dots, x_{n}$ coincide. Hence, we have
\begin{equation*}
	A_{\bm{\Phi}}=\Phi_{0}+\sum_{n=1}^{N}\sum_{\substack{x_{1},\dots,x_{n}\in\mfrak{X}\\ \mrm{distinct}}}\Phi_{n}(x_{1},\dots,x_{n})\prod_{i=1}^{n}a_{x_{i}}^{\ast}a_{x_{i}}.
\end{equation*}
We can also see that $\pi_{P_{0}}(a_{x}^{\ast}a_{x})e_{\omega}=\chi_{[x\in\omega]}e_{\omega}$, $x\in\mfrak{X}$, $\omega\in\Omega^{\circ}$. Therefore, the desired result is obtained.
\end{proof}

The eigenvalues of $\pi_{P_{0}} (A_{\bm{\Phi}} )$ are non-negative from the assumption implying that $A_{\bm{\Phi}}$ is a positive element, and therefore, the system of functions $\{\rho_{n}\}_{n=1}^{\infty}$ fulfills the positivity conditions.
Now, the proof is complete.

